\documentclass[twoside,11pt]{article}

\usepackage[preprint]{jmlr2e}
\usepackage{amsmath}
\usepackage{array}
\usepackage{booktabs}
\usepackage{lastpage}
\usepackage{microtype}
\usepackage{tabularx}

\graphicspath{%
  {sections/methodology/scripts/generated/}%
  {sections/results/scripts/generated/}%
}

\hypersetup{
  hidelinks,
  pdftitle={VOCODE: A Systematic Evaluation of Quality, Efficiency, and Deployment Trade-offs in Neural Vocoders},
  pdfauthor={Rahul Sawhney},
  pdfsubject={Neural vocoder quality, efficiency, and deployment optimization},
  pdfkeywords={neural vocoder, speech synthesis, model efficiency, real-time factor, deployment optimization}
}

\jmlrheading
  {1}
  {2026}
  {1-\pageref{LastPage}}
  {7/26}
  {}
  {vocode-course-project}
  {Rahul Sawhney}
\ShortHeadings
  {Quality, Efficiency, and Deployment Trade-offs in Neural Vocoders}
  {Sawhney}
\firstpageno{1}

\begin{document}

\title{VOCODE: A Systematic Evaluation of Quality, Efficiency, and
Deployment Trade-offs in Neural Vocoders}

\author{\name Rahul Sawhney \email rahul.sawhney@fau.de \\
  \addr Department Artificial Intelligence in Biomedical Engineering\\
  Friedrich-Alexander-Universit\"at Erlangen-N\"urnberg\\
  N\"urnberger Stra\ss e 74, 91052 Erlangen, Germany}

\maketitle

\begin{abstract}%
Neural vocoders differ sharply in quality, speed, and footprint, yet
published results resist comparison: corpus, training budget,
implementation, evaluator, and hardware vary across papers.  VOCODE
benchmarks twelve vocoders spanning autoregressive, GAN-based,
iSTFT-based, flow-matching, and attention-augmented designs, trained from
scratch on LJSpeech with architecture-native objectives and identical
data splits.  Evaluation is unified: ten quality metrics (PESQ, STOI,
UTMOS, spectral and pitch distances) plus real-time factor, latency,
parameters, and size.  Deployment interventions yield 69 transformed groups paired with 22
same-hardware baselines (273 executions), each judged under an explicit
decision rule: 19 beneficial, 32 harmful, 18 inconclusive.  No vocoder is Pareto-optimal across quality, latency, and
footprint: the speed-quality frontier is Vocos, RFWave, and RNDVoC.  Gains
were architecture-dependent: FP16 halved storage at preserved PESQ in
eight of ten groups, quality-preserving speedups reached $2.07\times$
(ONNX) and $1.40\times$ (compilation), and RFWave at eight ODE steps
gained
$1.20\times$ within 0.05 PESQ.  Masked pruning was harmful in 18 of 19;
static INT8 compressed models near-fourfold at material quality loss.  At
moderate quality, parameter, and storage budgets, selection chose
FP16-cast Vocos on GPU and dynamic-INT8 Vocos on CPU.  With an adaptive
evaluation set, unreplicated training, and no listening test, VOCODE
supports conditional, auditable comparisons, not a causal architecture
ranking.
\end{abstract}

\begin{keywords}
neural vocoder, speech synthesis, model efficiency, real-time factor,
quantization, pruning, Pareto analysis, reproducibility
\end{keywords}

\section{Introduction}
\label{sec:introduction}

A neural vocoder reconstructs speech from an acoustic representation.
Autoregressive systems generate samples sequentially
\citep{oord2016wavenet}, whereas adversarial, inverse short-time Fourier
transform (iSTFT), and flow models move learned computation toward parallel
waveform synthesis, frame-level spectral prediction, or a few solver steps
\citep{kong2020hifigan,siuzdak2024vocos,liu2025rfwave}.  These computational
graphs expose different bottlenecks, so parameter count, serial dependency,
memory traffic, transform overhead, and solver evaluations need not track
one another.

Published quality-efficiency rankings also mix corpora, preprocessing,
training budgets, checkpoint rules, implementations, evaluators, and
hardware.  Even within one implementation, parameter count, serialized
bytes, multiply-accumulate counts, and real-time factor (RTF) answer
different questions.  Deployment transformations widen the gap between a
method name and the artifact that executes: compilation can break a graph,
quantization can cover only selected operators, and dense pruning can retain
neither sparse storage nor sparse kernels
\citep{han2015pruning,jacob2018quantization,ansel2024pytorch2}.

This work benchmarks twelve neural-vocoder configurations trained from
scratch under identical data splits and one evaluation implementation.  The
cohort comprises APNet2, BigVGAN-base, FreeV, HiFi-GAN V1, V2, V3, LPCNet,
MelGAN, RFWave, RNDVoC, Vocos, and VocosFormer.  Architecture-native
conditioning, objectives, and training budgets are retained.  Each configuration
contributes one training trajectory from recorded seed 1234 and one retained
checkpoint.  Inference seeds 42, 43, and 44 re-execute that frozen state
rather than retrain; the single exception is static ONNX INT8, which
rebuilds calibration state on each execution.

Three research questions structure the study.  \emph{RQ1}: under identical
data splits and one evaluation implementation, but
architecture-native objectives and budgets, how do the twelve
configurations trade objective quality, inference latency,
parameter count, and deployable storage?  \emph{RQ2}: relative to
same-hardware baselines, how do executed deployment transformations shift
objective quality, warm real-time factor, and serialized artifact size,
and how strongly do those effects depend on architecture and execution
backend?  \emph{RQ3}: which measured deployment artifact is selected under
explicit hardware-specific constraints on quality, parameter count,
and deployable storage?  The three are dependent analyses of one
experiment: transformed variants inherit the trained checkpoints, selection
consumes their measurements, and unsupported transformations are excluded
on record.

The scope is deliberately conditional.  LJSpeech contains one speaker.  The
final 525 utterances form an adaptive evaluation set: their PESQ values
informed checkpoint decisions and one deterministic repair of RFWave's
inference sampler.
Objective measures are proxies, not listener judgments
\citep{rix2001pesq,taal2011stoi,saeki2022utmos}.  Population-level or causal claims are precluded by one training trajectory
per configuration, unequal architecture-native budgets, independently
scheduled training jobs, and the absence of listening or multi-speaker
experiments.

\section{Literature Review}
\label{sec:literature-review}

The literature separates into generation mechanisms, deployment
transformations, and comparative evaluation; the unqualified label
\emph{vocoder efficiency} conflates model structure, executable
representation, and measurement protocol.

\subsection{Generation Mechanisms}

WaveNet factorizes waveform likelihood sample by sample
\citep{oord2016wavenet}; LPCNet reduces learned synthesis through linear
prediction but retains a recurrent excitation generator
\citep{valin2019lpcnet}.  Adversarial vocoders instead use parallel
generators.  MelGAN combines multi-scale discrimination with feature matching
\citep{kumar2019melgan}; HiFi-GAN adds multi-period discrimination and
quality-footprint variants \citep{kong2020hifigan}; BigVGAN introduces
periodic activations and anti-aliased representations
\citep{lee2023bigvgan}.  These mechanisms differ in critical-path
parallelism, so capacity alone cannot predict latency.

iSTFT vocoders perform learned prediction near the acoustic-frame rate and
delegate waveform reconstruction to overlap-add; iSTFTNet established the
pattern by replacing a time-domain generator's final upsampling stages
with the inverse transform \citep{kaneko2022istftnet}.  APNet2 predicts amplitude
and phase spectra \citep{du2023apnet2}; Vocos predicts complex coefficients
with a constant-resolution ConvNeXt stack \citep{siuzdak2024vocos}; FreeV
uses a pseudo-inverse mel prior \citep{lv2024freev}; and RNDVoC separates
analytical range and learned null-space components \citep{li2025rndvoc}.
RFWave instead integrates a rectified flow in a multi-band spectral
representation \citep{liu2025rfwave}, making solver evaluations an explicit
latency-quality control.

VocosFormer is not a published architecture: the project inserts
residual-convolution and self-attention components into a capacity-matched
Vocos chassis, adopting them from the WavTokenizer decoder lineage
\citep{ji2025wavtokenizer}; the codec is cited for component provenance
only.

\subsection{Deployment Transformations}

Pruning changes both numerical weights and, only when supported, their
storage and execution representation.  Classical train-prune-retrain work
couples sparsity to compressed storage \citep{han2015pruning};
dense zero masks establish tolerance to weight removal but not
acceleration.
Quantization likewise specifies weights, activations, scales, zero-points,
covered operators, and backend \citep{jacob2018quantization}.  CPU dynamic
quantization, ONNX QDQ graphs, and GPU weight-only quantization are therefore
distinct treatments despite sharing an integer bit width.  Compilation
changes the captured executable graph; graph breaks, specialization mode,
operator support, and hardware determine whether TorchDynamo/TorchInductor
produces a gain \citep{ansel2024pytorch2}.  The vocoder literature itself
pursues efficiency through architecture and training: WaveRNN prunes its
recurrent weights during training to reach real-time mobile-CPU synthesis
\citep{kalchbrenner2018wavernn}, and WaveFM distills a flow-matching
vocoder toward single-step synthesis \citep{luo2025wavefm}; the
transformations above instead operate on fixed trained checkpoints.
Reproducible deployment
evidence must therefore state transformation scope, backend, operator
coverage, and the paired pre-transformation measurement, not transformation
names alone.

\subsection{Comparative Evaluation}

PESQ models perceived speech distortion \citep{rix2001pesq}, STOI estimates
intelligibility \citep{taal2011stoi}, and UTMOS predicts mean opinion scores
\citep{saeki2022utmos}; agreement among these proxies is still not a listening
test.  Shared-protocol vocoder comparisons improve internal validity while
remaining sensitive to implementation and conditioning choices
\citep{govalkar2019comparison}, and VocBench trains and evaluates six
vocoders in one shared environment with objective and subjective scoring
\citep{albadawy2022vocbench}.  Deployment-focused work shows that
batch speed does not establish low-latency streaming performance because
causality, chunk size, lookahead, loading overhead, and hardware can reverse
rankings \citep{yoneyama2025comparative}.  VOCODE does not claim streaming
evaluation or listener scoring.  Relative to VocBench it contributes twelve configurations trained from
scratch under recorded architecture-specific budgets, explicit checkpoint
provenance, inclusion and exclusion registers for executable deployment
artifacts, paired before-and-after measurements on one hardware profile,
and constraint-aware selection.

\section{Methodology}
\label{sec:methodology}

We benchmark twelve neural vocoders under one common data, evaluation, and
measurement protocol.  A \emph{Project-Trained Configuration} is one architecture recipe
fitted from random initialization under the project budget; its selected
state is the \emph{Retained Project Checkpoint}, which every measurement
inherits.  The protocol fixes corpus membership, evaluator, and measurement
procedure while each configuration keeps its native conditioning, objective,
and optimization recipe: a structured common testbed, not an equal-compute
causal benchmark.  Throughout, a deployment artifact $a=(m,t,h)$ applies
transformation $t$ to the checkpoint $\theta_m^\star$ of configuration
$m\in\mathcal M$ on requested \emph{hardware lane}
$h\in\{\mathrm{B200},\mathrm{CPU}\}$; the
identity transformation defines the paired same-lane control $c(a)$, and
inference seeds $s\in\{42,43,44\}$ re-execute one selected state without
resampling training (single exception: static ONNX INT8,
Section~\ref{sec:deployment-methods}).

\subsection{Data Set and Preprocessing}
\label{sec:data-preprocessing}

All experiments use LJSpeech~1.1, 13,100 English utterances from one speaker
\citep{ito2017ljspeech}.  Identifier-sorted disjoint blocks fix 12,475
training, 100 validation, and 525 evaluation items without a random seed;
identifier order preserves chapter structure, so the split is fixed
membership rather than a random utterance sample.
Training loaders shuffle and crop fixed-length segments,
whereas validation and evaluation consume complete utterances.  The
evaluation suffix is \emph{adaptive}, not untouched: its PESQ values informed
stop and extension decisions, and one deterministic repair of RFWave's
inference sampler.

Every configuration retains its native conditioning interface: most consume
80 Slaney-scale mel bands at 22.05~kHz; BigVGAN-base uses 100 normalized
bands, and Vocos, VocosFormer, and RFWave use 100 unnormalized HTK-scale
bands at 24~kHz; all mel pipelines share a 1,024-point Hann transform with
a 256-sample hop, while LPCNet consumes 18 Bark-frequency cepstra plus pitch and
correlation features at 16~kHz.  Training-only transformations are peak
normalization to 0.95 (HiFi-GAN V1, V2, and V3), uniform random peak gain in
$[-6,-1]$~dB (Vocos, VocosFormer), and LPCNet excitation noise with standard
deviation 0.3; evaluation waveforms receive none.

\subsection{Model Architectures}
\label{sec:model-architectures}

The cohort spans five generation mechanisms: time-domain adversarial
upsamplers (MelGAN, HiFi-GAN V1, V2, V3, BigVGAN-base), adversarially
trained spectral models with inverse-STFT heads (APNet2, FreeV, RNDVoC,
Vocos), hybrid attention (VocosFormer), sample-level autoregression
(LPCNet), and rectified-flow synthesis (RFWave);
Table~\ref{tab:benchmark-cohort} reports the executed configuration surface,
with the cited architectural definitions in the literature review.

\begin{table}[t]
  \centering
  \footnotesize
  \renewcommand{\arraystretch}{1.02}
  \setlength{\tabcolsep}{4pt}
  \InputIfFileExists{sections/methodology/scripts/generated/cohort_table.tex}{}{%
    \begin{tabular*}{\linewidth}{@{\extracolsep{\fill}}lllrrrcrr@{}}
  \toprule
  Configuration & Family & Conditioning & Params & Segment & Batch & Prec. & Updates & State \\
  \midrule
  APNet2 & iSTFT GAN & 80 mel / 22.05 & 31.43 & 8,192 & 32 & BF16 & 124.8 & 120k \\
  BigVGAN-base & Waveform GAN & 100 mel / 24 & 14.03 & 8,192 & 32 & BF16 & 124.8 & 120k \\
  FreeV & iSTFT GAN & 80 mel / 22.05 & 18.22 & 8,192 & 32 & BF16 & 124.8 & 120k \\
  HiFi-GAN V1 & Waveform GAN & 80 mel / 22.05 & 13.94 & 8,192 & 32 & BF16 & 254.8 & 245k \\
  HiFi-GAN V2 & Waveform GAN & 80 mel / 22.05 & 0.93 & 8,192 & 32 & BF16 & 124.8 & 120k \\
  HiFi-GAN V3 & Waveform GAN & 80 mel / 22.05 & 1.46 & 8,192 & 32 & BF16 & 254.8 & 245k \\
  LPCNet & Autoregressive & 20 feat / 16 & 1.23 & 2,400 & 128 & FP32 & 31.3 & 30k \\
  MelGAN & Waveform GAN & 80 mel / 22.05 & 4.27 & 16,000 & 32 & BF16 & 124.8 & 120k \\
  RNDVoC & iSTFT GAN & 80 mel / 22.05 & 3.57 & 16,384 & 32 & BF16 & 124.8 & 120k \\
  RFWave & Rectified flow & 100 mel / 24 & 18.14 & 32,512 & 64 & BF16 & 125.0 & 120k \\
  Vocos & iSTFT GAN & 100 mel / 24 & 13.53 & 16,384 & 32 & BF16 & 124.8 & 120k \\
  VocosFormer & iSTFT GAN + attn & 100 mel / 24 & 13.78 & 16,384 & 32 & BF16 & 124.8 & 120k \\
  \bottomrule
\end{tabular*}
}
  \caption{Project-Trained Configurations.  Conditioning: input features and
  rate (kHz); Params: deployable parameters (millions); Segment: training
  crop (samples); Updates: executed budget (thousands); State: the Retained
  Project Checkpoint, the terminal durable state selected by the
  registered gate.}
  \label{tab:benchmark-cohort}
\end{table}

VocosFormer, the project-defined exception, inserts a residual-convolution,
self-attention, residual-convolution block from the WavTokenizer decoder
lineage into a shallower Vocos chassis under matched losses, data, and
training quantum; its 13.78 million parameters sit $1.83\%$ above Vocos,
and one training trajectory supports only configuration-level comparison.

\subsection{Training and Checkpoint Selection}
\label{sec:training-checkpoint}

We fitted each configuration through one recorded-seed (1234) training
trajectory under its native objective and a finite registered budget.  All
models are PyTorch modules executed through the project-controlled
SyntheticMind harness, which standardizes fitting, validation, metric
reduction, and checkpoint restore; the VOCODE library
supplies data pipelines, objectives, synthesis paths, evaluators, and
deployment transformations, and each configuration, variant, lane, and
execution runs as one isolated job.
Eleven runs used a 48-GB NVIDIA L40S and RFWave an 80-GB NVIDIA A100;
recorded training sums to 75.31 GPU-hours, a lower bound.
Architecture-native rather than equalized budgets make the retained states
controlled project realizations, not attainable quality ceilings.

Objectives are architecture-native.  MelGAN optimizes least-squares
adversarial and feature-matching terms alone; HiFi-GAN V1, V2, and V3 and
BigVGAN-base combine the same terms over their discriminator ensembles
with weighted mel reconstruction, and Vocos and VocosFormer apply the
hinge form of that composition.  APNet2, FreeV, and
RNDVoC add log-amplitude, anti-wrapping phase, and STFT-consistency terms
with real and imaginary reconstruction to mel, hinge adversarial, and
feature-matching supervision.  LPCNet minimizes teacher-forced cross
entropy over its quantized mu-law excitation; RFWave regresses conditional
flow velocities in the waveform domain with auxiliary magnitude and
band-overlap terms.  Objectives were deliberately not standardized, so raw
training and validation losses are not comparable across families and
never ranked configurations.  Ten configurations optimized with AdamW and
two (MelGAN, LPCNet) with Adam; initial learning rates were
$2{\times}10^{-4}$ (APNet2, FreeV, HiFi-GAN V1, V2, and V3, RNDVoC,
RFWave), $1{\times}10^{-4}$ (BigVGAN-base, MelGAN), $5{\times}10^{-4}$
(Vocos, VocosFormer), and $1{\times}10^{-3}$ (LPCNet).  LPCNet additionally
trains under scheduled GRU sparsification from 2k to 20k updates, the
HiFi-GAN V1 and V3 budgets include 4,800 replayed updates after a durable
resume, and exact coefficients and schedules remain recoverable from the
resolved run records and code.

Checkpoint selection follows a registered gate: validation loss informed
continuation, durable checkpoints were written every 5,000 updates, a
sub-0.02 PESQ gain across gates indicated a plateau, and every run kept a
registered ceiling.  Where a commensurate published or released-checkpoint
PESQ anchor existed, the registered quality target was 93\% of that
anchor; RFWave used a registered PESQ floor of 3.55, and LPCNet and
VocosFormer had no commensurate numeric anchor.  For each configuration, evaluation
used the terminal durable checkpoint admitted by this gate, the Retained
Project Checkpoint $\theta_m^\star$ of Table~\ref{tab:benchmark-cohort},
not a retrospectively substituted validation minimum; the
validation-best, global-minimum, and evaluated states remain in the
resolved records.  The accepted runs predate the repair that seeds model
construction and workers deterministically, so 1234 is a recorded rather
than reconstructible seed; retained states are immutable, and bitwise
regeneration is not claimed.

\subsection{Deployment Optimization Methods}
\label{sec:deployment-methods}

We transform each Retained Project Checkpoint with the deployment variants
of Figure~\ref{fig:deployment-admission}, spanning TorchInductor
compilation, FP16 weight casting, INT8 quantization, ONNX deployment,
magnitude pruning at 30/50/70\% with one recovery arm, and RFWave
solver-step reduction.  Admission requires an executable transformed path,
nonzero operator or parameter coverage, all 525 utterances under three
executions, and a same-profile pre-transformation measurement; unsupported
cells fail before execution expansion.  The executed surface contains 22
pre-transformation controls and 69 transformed groups (91 groups, 273
executions); LPCNet completed no admissible execution on either lane,
leaving eleven transformed configurations, and 29 exclusion records
document 88 rejected attempts.

\begin{figure}[t]
  \centering
  \includegraphics[width=0.88\linewidth]{deployment_admission_matrix}
  \caption{Deployment admission surface.  Filled: admitted three-execution
  groups (91, incl.\ 22 same-lane controls); crosses: the ten excluded
  attempts inside the plotted vocabulary (invalid recovery or control
  states; incomplete LPCNet executions), with the 29-row exclusion register
  remaining authoritative.  Compile RO = reduce-overhead; INT8 WO =
  weight-only; Recovery = 13-epoch recovery of the 50\%-pruned state; ODE
  $k$ = $k$ Euler steps.}
  \label{fig:deployment-admission}
\end{figure}

Static ONNX INT8 uses ONNX
Runtime~1.24.4 on its CPU provider with eight intra-op threads, QDQ format,
per-channel QInt8 weights, QUInt8 activations, and coverage of Conv,
ConvTranspose, MatMul, and Gemm.  Each execution reconstructs its quantized
graph, calibrating on 48 mel batches from the execution-seeded shuffled
training loader, so these executions combine calibration-set and runtime
variation rather than repeating one fixed artifact.  Dynamic and weight-only
INT8 retain their distinct PyTorch and torchao semantics; an unsupported
INT4 probe contributes no result.  Pruning is magnitude-based with dense
masks, testing weight removal but not sparse storage or kernels: Vocos,
FreeV, VocosFormer, and RFWave use $L_2$-ranked structured output-channel
pruning over backbone Linear projections (spectral heads excluded); the rest
use global $L_1$ unstructured pruning over convolutional, Linear, and
recurrent weights.  Selected 50\% states received 13 recovery epochs at one
tenth learning rate, masks held during fine-tuning and zeros baked into the
evaluated checkpoint; absent an admitted dense-continuation control,
recovered rows support descriptive rather than causal claims.  RFWave step
reduction shrinks its ten-step Euler baseline to eight, four, or two steps.

\subsection{Evaluation Metrics and Analysis Protocols}
\label{sec:evaluation-protocols}

\paragraph{Quality proxies.}
Wideband PESQ, withdrawn by ITU-T in January 2024 in favor of the P.863
family \citep{itut2024p862} yet still the dominant quality proxy in
the vocoder literature, remains the primary intrusive measure
\citep{rix2001pesq};
STOI \citep{taal2011stoi} and UTMOS \citep{saeki2022utmos} test dependence on
the proxy choice.  Reference and synthesized waveforms are truncated to
their common unpadded length (\emph{true-length} evaluation) before
resampling; PESQ operates wideband at
16~kHz, classic non-extended STOI at 10~kHz, and UTMOS (SpeechMOS UTMOS22
strong~1.2.0) scores the synthesized waveform alone at 16~kHz.  Mel $L_1$, multi-resolution STFT, mel-cepstral,
log-spectrum, pitch, periodicity, and voicing errors are diagnostics only;
architecture-native or alignment-sensitive values are not cross-model
endpoints, and without a blinded listening study every quality claim remains
a proxy claim.  For utterance $i$ in the fixed 525-item evaluation set
$\mathcal I$, with conditioning $x_i$ and reference $y_i$, artifact $a$
under execution $s$ synthesizes $\widehat y_{a,s,i}=f_{a,s}(x_i)$, and
proxy $q$ is aggregated as
\begin{equation}
  \widehat Q_a^{(q)}
  =\frac{1}{|\mathcal I|}\sum_{i\in\mathcal I}
   \frac{1}{|\mathcal S|}\sum_{s\in\mathcal S}
   q\bigl(y_i,\widehat y_{a,s,i}\bigr),
  \qquad \mathcal S=\{42,43,44\},
  \label{eq:quality-estimator}
\end{equation}
where the reference-free UTMOS ignores $y_i$; execution values are averaged
within each utterance before any paired contrast, and for static ONNX INT8
the execution index also carries calibration-state variation
(Section~\ref{sec:deployment-methods}).

\paragraph{Timing and footprint protocol.}
Beyond quality, endpoints include the continuous process-RSS high-water
mark, one-time ONNX session-construction time, and complexity as
parameters, serialized bytes, and MACs from one traced second of audio;
jobs were independently
scheduled, and paired contrasts match the requested resource profile,
never the physical host.  After five warm-up batches,
pre-transformation evaluation measures one complete pass at batch one on
an NVIDIA B200.  The CPU column of Table~\ref{tab:baseline-surface} is the
retained comparative record, three executions each timing one complete
post-warm-up pass on an 8-core CPU host, and is never cross-compared with
deployment-phase RTFs; LPCNet's 27-item timing record is excluded from
strict speed frontiers.  Deployment evaluation measures three synchronized
calls per utterance for the 520 post-warm-up utterances on B200 or CPU,
each timed call
following one untimed prediction of the same input, so the quantities
describe warm repeated synthesis only.  For artifact $a$, execution $s$,
timed items $\mathcal I_a\subseteq\mathcal I$, synchronized time
$\tau_{a,s,i,r}$, duration $d_i$, and $J_a\in\{1,3\}$ calls, warm RTF is
\begin{equation}
  \bar\tau_{a,s,i}=J_a^{-1}\sum_{r=1}^{J_a}\tau_{a,s,i,r},\qquad
  R_{a,s}=\frac{\sum_{i\in\mathcal I_a}\bar\tau_{a,s,i}}
                 {\sum_{i\in\mathcal I_a}d_i},
  \label{eq:rtf}
\end{equation}
and the group-level $\widehat R_a$ is the mean of $R_{a,s}$ over
executions.  Reported p50/p95 are per-execution percentiles of the 520
per-utterance three-repetition means $\bar\tau_{a,s,i}$, averaged over
executions, not of the 1,560 individual calls, whose raw durations were
not retained.  B200 jobs requested 16 CPUs and 64~GiB, CPU
jobs 8 CPUs and 16~GiB; the Python~3.14.2, PyTorch~2.13.0+cu130,
torchaudio~2.11.0, torchao~0.17.0, ONNX~1.20.1, and ONNX Runtime~1.24.4 stack
is retained, whereas host-identifying fields are irrecoverable.

\paragraph{Paired contrasts and decision rule.}
For quality proxy $q$ and transformed artifact $a$ with paired control
$c(a)$, quality estimates $\widehat Q$ from
Eq.~\eqref{eq:quality-estimator}, warm RTF $\widehat R$ from
Eq.~\eqref{eq:rtf}, and serialized deployable size $D$, the continuous
effects are
\begin{equation}
  \Delta Q_a^{(q)}=\widehat Q_a^{(q)}-\widehat Q_{c(a)}^{(q)},\qquad
  S_a=\widehat R_{c(a)}/\widehat R_a,\qquad
  Z_a=D_a/D_{c(a)}.
  \label{eq:deployment-effects}
\end{equation}
Quality deltas are utterance-paired, whereas speedups and size ratios are
same-profile descriptive contrasts because hosts are not paired.  Quality
contrasts pair the execution-averaged utterance values of
Eq.~\eqref{eq:quality-estimator} before a Student $t$ interval, with a
5,000-replicate percentile bootstrap as sensitivity analysis
\citep{student1908probableerror,efron1979bootstrap}.  These are fixed-set
descriptive intervals quantifying utterance dispersion over the 525 fixed
items under a working exchangeability assumption; they omit multiplicity
correction and variation from training, speaker, recovery seed, or host,
and for static INT8 do not separate calibration from runtime variation.
A post hoc ordered classifier provides an exploratory engineering summary
of Eq.~\eqref{eq:deployment-effects}: dense-masked pruning and recovery rows
are labeled first: harmful when $\Delta Q_a^{(\mathrm{PESQ})}<-0.10$,
otherwise inconclusive; remaining MelGAN rows are fixed as inconclusive
because their PESQ is floor-censored.  Every other row receives
\begin{equation}
  \ell(a)=
  \begin{cases}
    \text{harmful}, &
      \Delta Q_a^{(\mathrm{PESQ})}<-0.10
      \ \lor\ \bigl(S_a<0.90\wedge Z_a>0.90\bigr)\\[1pt]
    \text{beneficial}, &
      \Delta Q_a^{(\mathrm{PESQ})}\ge-0.05
      \ \wedge\ S_a\ge0.90
      \ \wedge\ \bigl(S_a\ge1.10\lor Z_a\le0.55\bigr)\\[1pt]
    \text{inconclusive}, & \text{otherwise.}
  \end{cases}
  \label{eq:decision-rule}
\end{equation}
Thresholding PESQ alone makes these labels PESQ-preserving statements, not
perceptual equivalence claims.  A $3^5=243$ grid tests label sensitivity by
varying each cutoff in Eq.~\eqref{eq:decision-rule} over three levels
(preservation $-0.075/-0.05/-0.025$; harm $-0.15/-0.10/-0.075$; speed floor
$0.85/0.90/0.95$; gain $1.05/1.10/1.20$; size $0.50/0.55/0.60$).

\paragraph{Frontier and constraint-aware selection.}
Under one quality proxy $q$ and one cost axis (RTF, parameter count, or
serialized size), artifact $a$ dominates $a'$ when
$\widehat Q_a^{(q)}\ge\widehat Q_{a'}^{(q)}$ at no greater cost, with at
least one inequality strict; each reported frontier is the nondominated
set, recomputed per proxy and axis.  For
hardware lane $h$, admitted artifacts $\mathcal A_h$, quality floor
$q_{\min}$, parameter ceiling $P_{\max}$, and size ceiling $D_{\max}$,
deployment selection solves
\begin{equation}
  a_h^\star\in\underset{a\in\mathcal A_h}{\arg\min}\ \widehat R_a
  \quad\text{s.t.}\quad
  \widehat Q_a^{(q)}\ge q_{\min},\quad P_a\le P_{\max},\quad
  D_a\le D_{\max},
  \label{eq:constrained-selection}
\end{equation}
so every recommendation is conditional on hardware, proxy, and constraints.

\section{Results}
\label{sec:results}

Results cover the twelve Project-Trained Configurations and the 91 admitted
model-variant groups of Figure~\ref{fig:deployment-admission}: 22 same-lane
controls and 69 transformed groups, measured over 273 executions.  All
quality values are the estimator of Eq.~\eqref{eq:quality-estimator} on the
fixed adaptive LJSpeech evaluation set, and every claim concerns these executed
states.  Except for static ONNX INT8
(Section~\ref{sec:deployment-methods}), the three executions per group
repeat inference from one frozen artifact, so their spread reflects
execution variability rather than statistical uncertainty.

\subsection{Baseline Quality, Efficiency, and Footprint}
\label{sec:baseline-results}

Addressing RQ1, Table~\ref{tab:baseline-surface} reports the complete
twelve-configuration outcome surface, including the floor-level negative
results, and
Figure~\ref{fig:baseline-landscape} places every configuration on the
quality-speed plane.  True-length PESQ ranked RNDVoC (3.895), RFWave
(3.785), and Vocos (3.614) as the three strongest states, and STOI produced
the same leading order; seven configurations reached true-length PESQ of at
least 3.30.  UTMOS disagreed with the intrusive proxies: HiFi-GAN V1 was
highest at 4.203, followed by Vocos (4.089) and BigVGAN-base (4.062); no
configuration was uniformly best across the three proxies.

\begin{table}[t]
  \centering
  \footnotesize
  \renewcommand{\arraystretch}{1.02}
  \setlength{\tabcolsep}{4pt}
  \InputIfFileExists{sections/results/scripts/generated/baseline_table.tex}{}{%
    \begin{tabular*}{\linewidth}{@{\extracolsep{\fill}}lrrrrrrr@{}}
  \toprule
  Configuration & Params & Size & PESQ & STOI & UTMOS & RTF B200 & RTF CPU$^{\mathsection}$ \\
  \midrule
  APNet2 & 31.43 & 119.88 & 2.515 & 0.930 & 3.091 & 0.0095 & 0.0191 \\
  BigVGAN-base & 14.03 & 53.51 & 3.320 & 0.974 & 4.062 & 0.1573 & 0.6985 \\
  FreeV & 18.22 & 69.66 & 3.386 & 0.969 & 3.789 & 0.0077 & 0.0157 \\
  HiFi-GAN V1 & 13.94 & 53.16 & 3.305 & 0.974 & 4.203 & 0.1197 & 0.1836 \\
  HiFi-GAN V2$^{\ddagger}$ & 0.93 & 3.54 & 2.400 & 0.942 & 3.907 & 0.1140 & 0.0507 \\
  HiFi-GAN V3$^{\ddagger}$ & 1.46 & 5.59 & 2.550 & 0.949 & 3.851 & 0.0580 & 0.0366 \\
  LPCNet & 1.23 & 4.70 & 1.052 & 0.527 & 1.235 & 11.36$^{*}$ & 9.2015 \\
  MelGAN & 4.27 & 16.27 & 1.043 & 0.418 & 1.736 & 0.0742 & 0.0456 \\
  RFWave$^{\dagger}$ & 18.14 & 69.21 & 3.785 & 0.981 & 3.514 & 0.0148 & 0.7491 \\
  RNDVoC$^{\dagger}$$^{\ddagger}$ & 3.57 & 14.94 & 3.895 & 0.987 & 4.007 & 0.0559 & 0.0838 \\
  Vocos$^{\dagger}$ & 13.53 & 51.62 & 3.614 & 0.979 & 4.089 & 0.0047 & 0.0098 \\
  VocosFormer & 13.78 & 52.57 & 3.512 & 0.976 & 3.901 & 0.0122 & 0.0105 \\
  \bottomrule
\end{tabular*}
}
  \caption{Baseline outcome surface for all twelve configurations: quality
  proxies (higher is better), parameters (millions), serialized size (MB),
  and warm real-time factors (lower is better).
  $^{\mathsection}$retained comparative CPU record
  (Section~\ref{sec:evaluation-protocols}), not a deployment-phase value;
  $^{\dagger}$B200 speed-quality frontier;
  $^{\ddagger}$parameter- and size-quality frontier;
  $^{*}$LPCNet timing (27 utterances, one execution) is excluded from
  strict frontiers.}
  \label{tab:baseline-surface}
\end{table}

\begin{figure}[t]
  \centering
  \includegraphics[width=\linewidth]{baseline_landscape}
  \caption{Twelve frozen baselines on the true-length PESQ versus B200
  real-time-factor plane (Eq.~\eqref{eq:rtf}).  The legend identifies each
  configuration, ordered by PESQ; marker shape encodes the generation
  family (squares: frame-level Fourier; circles: time-domain GAN;
  triangle: flow matching; diamond: hybrid attention; inverted triangle:
  autoregressive).  Dark-edged markers are frontier members, connected by
  the dashed nondominance chain; the dotted line marks real-time
  synthesis.  LPCNet's single-execution timing is shown but excluded from
  the strict frontier.  Each point is one row of the retained comparative
  record.}
  \label{fig:baseline-landscape}
\end{figure}

No checkpoint minimized all deployment costs.  Vocos had the lowest strict
B200 RTF (0.00473), RNDVoC attained the highest PESQ at 11.8 times that
RTF, and RFWave occupied the intermediate point (3.785 at 0.0148); the
strict speed frontier comprised Vocos, RFWave, and RNDVoC.  HiFi-GAN V2 was the smallest baseline (0.93 million
parameters, 3.54~MB) at PESQ 2.400; the parameter and size frontiers
comprised HiFi-GAN V2, HiFi-GAN V3, and RNDVoC.  Two configurations
produced floor-level negative results: MelGAN (PESQ 1.043 on both lanes)
is consistent with budget-limited non-convergence, and LPCNet (1.052)
lost free-running stability at 5.8\% of its reference recipe while its
PyTorch timing omits the author's optimized sparse C path; both remain
full members of the twelve-model cohort.

The matched VocosFormer state did not improve on its Vocos control: it added
247{,}040 parameters ($+1.83\%$), reduced true-length PESQ by 0.101 and STOI
by 0.0035, and increased B200 RTF $2.58\times$.  Under padded batch-16
evaluation the same checkpoint measured PESQ 2.480 against 3.512 at true
length, a padding-sensitivity failure consistent with contamination through
unmasked global attention rather than intrinsic synthesis collapse; this
one-trajectory comparison is a configuration-level observation, not an
average attention effect.

\subsection{Deployment Transformation Effects}
\label{sec:intervention-results}

Addressing RQ2, the decision rule of Eq.~\eqref{eq:decision-rule} labeled
the 69 transformed groups 19 beneficial, 18 inconclusive, and 32 harmful.
Table~\ref{tab:family-effects} reports the seven intervention families and
their fourteen variants with per-variant label composition, median
same-lane effects, and $\Delta$PESQ ranges; the 29-row exclusion
register documents the 88 attempts producing no admitted group.  FP16
weight casting was the most consistently favorable family: eight of
ten groups were beneficial at half serialized size, with RFWave and
HiFi-GAN V3 reaching $1.955\times$ and $1.874\times$ speedups at preserved
PESQ.

\begin{table}[t]
  \centering
  \footnotesize
  \renewcommand{\arraystretch}{1.02}
  \InputIfFileExists{sections/results/scripts/generated/variant_table.tex}{}{%
    \begin{tabular*}{\linewidth}{@{\extracolsep{\fill}}lrcrcrr@{}}
  \toprule
  Variant & $n$ & B/I/H & Med.\ $\Delta$PESQ & $\Delta$PESQ range & Med.\ $S_a$ & Med.\ $Z_a$ \\
  \midrule
  \multicolumn{7}{@{}l}{\emph{Compilation}} \\
  \quad Default mode & 11 & 4/5/2 & $-0.0000$ & $[-0.000, +0.000]$ & 1.07 & 1.00 \\
  \quad Reduce-overhead mode & 4 & 1/0/3 & $-0.0000$ & $[-0.000, +0.000]$ & 0.40 & 1.00 \\
  \addlinespace[2pt]
  \multicolumn{7}{@{}l}{\emph{FP16}} \\
  \quad Weight cast & 10 & 8/2/0 & $-0.0002$ & $[-0.002, +0.001]$ & 1.08 & 0.50 \\
  \addlinespace[2pt]
  \multicolumn{7}{@{}l}{\emph{INT8}} \\
  \quad Dynamic & 5 & 2/1/2 & $-0.0469$ & $[-0.648, +0.004]$ & 1.30 & 0.40 \\
  \quad Weight-only & 4 & 2/2/0 & $-0.0010$ & $[-0.003, +0.002]$ & 0.92 & 0.41 \\
  \addlinespace[2pt]
  \multicolumn{7}{@{}l}{\emph{ONNX}} \\
  \quad FP32 export & 4 & 1/2/1 & $-0.0000$ & $[-0.000, +0.000]$ & 0.90 & 1.00 \\
  \quad INT8 static & 4 & 0/1/3 & $-0.7701$ & $[-1.329, -0.004]$ & 0.82 & 0.26 \\
  \addlinespace[2pt]
  \multicolumn{7}{@{}l}{\emph{Pruning}} \\
  \quad 30\% sparsity & 4 & 0/0/4 & $-1.0259$ & $[-2.740, -0.148]$ & 0.95 & 1.00 \\
  \quad 50\% sparsity & 11 & 0/1/10 & $-1.0632$ & $[-2.719, +0.000]$ & 1.00 & 1.00 \\
  \quad 70\% sparsity & 4 & 0/0/4 & $-2.2722$ & $[-2.712, -2.014]$ & 1.03 & 1.00 \\
  \addlinespace[2pt]
  \multicolumn{7}{@{}l}{\emph{Recovery}} \\
  \quad 50\% + recovery & 5 & 0/4/1 & $-0.0070$ & $[-0.475, +0.017]$ & 1.11 & 1.00 \\
  \addlinespace[2pt]
  \multicolumn{7}{@{}l}{\emph{ODE steps}} \\
  \quad 8 steps & 1 & 1/0/0 & $-0.0456$ & n/a & 1.20 & 1.00 \\
  \quad 4 steps & 1 & 0/0/1 & $-0.3689$ & n/a & 2.20 & 1.00 \\
  \quad 2 steps & 1 & 0/0/1 & $-1.6363$ & n/a & 4.51 & 1.00 \\
  \bottomrule
\end{tabular*}
}
  \caption{Deployment intervention outcomes over the 69 admitted groups,
  by family and variant: admitted count, beneficial/inconclusive/harmful
  composition under Eq.~\eqref{eq:decision-rule}, median same-lane
  $\Delta$PESQ, speedup $S_a$, and size ratio $Z_a$
  (Eq.~\eqref{eq:deployment-effects}), and the $\Delta$PESQ range across
  the variant's admitted groups; single-group variants report their one
  contrast.}
  \label{tab:family-effects}
\end{table}

Compilation was mode specific: default-mode compilation preserved PESQ in
all eleven groups at speedups from $0.14\times$ to $1.403\times$
(HiFi-GAN V3), whereas the reduce-overhead mode was harmful for three of
four groups (speedups 0.34 to 0.41) and beneficial only for RFWave
($1.361\times$).  ONNX FP32
accelerated HiFi-GAN V2 by $2.071\times$ at $\Delta$PESQ $-0.00013$
(fixed-set 95\% interval $[-0.00036, 0.00010]$) yet split 1/2/1 (B/I/H)
overall.  Dynamic INT8 diverged: beneficial for APNet2 and
Vocos, harmful for FreeV ($-0.123$) and RFWave ($-0.648$).  Weight-only
INT8 preserved PESQ in all four groups but slowed FreeV ($0.851\times$) and
VocosFormer ($0.617\times$).

Static ONNX INT8 was harmful in three of four groups with median PESQ loss
0.770 at median size ratio 0.265.  Because each execution recalibrates on
its own shuffled loader batches, its per-execution PESQ disperses where all
other variants repeat exactly: HiFi-GAN V1 measured 2.003, 2.053, and
1.910 against its execution-invariant ONNX FP32 value of 3.317.  Dense magnitude
masks were harmful in 4/4, 10/11, and 4/4 groups at 30/50/70\% sparsity
(median $\Delta$PESQ $-1.03$, $-1.06$, and $-2.27$); the single non-harmful
group is floor-censored MelGAN.  Recovery returned four of five admitted
50\% groups to within 0.05 PESQ of their controls (BigVGAN-base moved
from $-1.149$ to $+0.017$) but left Vocos harmful ($-0.475$); lacking
the dense-continuation control, the recovery arm supports no causal claim.  RFWave
solver reduction traded monotonically: speedups of 1.201, 2.200, and 4.509
for PESQ losses of 0.046, 0.369, and 1.636 at 8, 4, and 2 steps.

Deployment-phase RTFs follow Eq.~\eqref{eq:rtf} with $J_a{=}3$, not
comparable with Table~\ref{tab:baseline-surface}'s single-pass B200 or
comparative CPU baselines (Section~\ref{sec:evaluation-protocols}).  Across the 273 executions, p50 spanned
1.21 to 3{,}909.35~ms, p95 1.96 to 6{,}138.02~ms, process RSS 6{,}754 to
16{,}991~MiB, and the 24 ONNX session builds spanned 0.0116 to 0.1327~s.  The
deployment costs need not move together: the HiFi-GAN V2 ONNX FP32
artifact halves median synthesis time on CPU while raising process RSS
above its control.  Of the 69 utterance-paired PESQ contrasts, 34 fixed-set intervals were
entirely negative, two entirely positive, and 33 crossed zero.  Despite
non-normality (113 contrasts with absolute skewness above one, 181 with
excess kurtosis above three), bootstrap and Student $t$ signs agreed on
482 of 483 contrasts, an internal agreement diagnostic, not external
validity.

\subsection{Sensitivity of the Decision Taxonomy and Constraint-Aware
Selection}
\label{sec:selection-results}

Across the $3^5=243$ threshold profiles of
Section~\ref{sec:evaluation-protocols}, the beneficial,
harmful, and inconclusive totals ranged over 12 to 24, 29 to 34, and 12 to
27; median and minimum agreement with the default taxonomy were 91.3\%
and 82.6\%.  Stability stratifies by rule structure: 27 of the 28 rows
governed by the dense-mask and MelGAN overrides never changed label
(96.4\%), whereas 25 of the 41 threshold-governed rows never changed
(61.0\%), so the aggregate stability (52 of 69 rows unchanged) partly
reflects deterministic rule structure rather than threshold insensitivity.

Frontier membership depended on the proxy: STOI reproduced the
PESQ frontiers of Section~\ref{sec:baseline-results}, while UTMOS yielded
Vocos then HiFi-GAN V1 against B200 RTF and, against parameters, HiFi-GAN
V2, RNDVoC, Vocos, and HiFi-GAN V1, so proxy choice alters which states appear
efficient.

\begin{table}[t]
  \centering
  \footnotesize
  \renewcommand{\arraystretch}{1.02}
  \InputIfFileExists{sections/results/scripts/generated/selector_table.tex}{}{%
    \begin{tabular*}{\linewidth}{@{\extracolsep{\fill}}llrlrr@{}}
  \toprule
  Constraints $(q_{\min}, P_{\max}, D_{\max})$ & Lane & Feasible & Selected artifact $a_h^\star$ & RTF & PESQ \\
  \midrule
  (3.30, 20M, 70\,MB) & B200 & 28 & FP16-cast Vocos & 0.0003 & 3.611 \\
  (3.30, 20M, 70\,MB) & CPU & 8 & Dynamic-INT8 Vocos & 0.0055 & 3.563 \\
  (3.60, 15M, 55\,MB) & B200 & 10 & FP16-cast Vocos & 0.0003 & 3.611 \\
  (3.60, 15M, 55\,MB) & CPU & 2 & Untransformed Vocos & 0.0067 & 3.610 \\
  \bottomrule
\end{tabular*}
}
  \caption{Constraint-aware selection (Eq.~\eqref{eq:constrained-selection})
  over all admitted artifacts per lane, with same-lane controls:
  feasible-set size and RTF-minimizing artifact per constraint vector
  (deployment-phase PESQ floor, parameter ceiling, size ceiling).  RTF
  values follow the deployment-phase three-repetition estimand (not
  comparable with Table~\ref{tab:baseline-surface}).}
  \label{tab:joint-selector}
\end{table}

Answering RQ3, Table~\ref{tab:joint-selector} instantiates the selector of
Eq.~\eqref{eq:constrained-selection} over all admitted artifacts of each
lane, using deployment-phase quality, RTF, deployable size, and logical
parameter counts consistently.  Under the moderate operating point
$(q_{\min}{=}3.30, P_{\max}{=}20\mathrm{M}, D_{\max}{=}70~\mathrm{MB})$,
28 B200 artifacts and eight CPU artifacts remain feasible, and the selector
chooses FP16-cast Vocos on B200 and dynamic-INT8 Vocos on CPU; tightening
to $(3.60, 15\mathrm{M}, 55~\mathrm{MB})$ leaves ten and two feasible
artifacts, selecting FP16-cast Vocos on B200 and the untransformed Vocos
control on CPU.  The feasible set remains nonempty on the B200 lane for
quality floors up to RNDVoC's FP16 artifact at 3.894; above that value no
admitted artifact qualifies.

\section{Discussion}
\label{sec:discussion}

RQ1 has no scalar answer.  RNDVoC's range-null
decomposition produced the strongest intrusive-proxy state at modest
parameter count, Vocos converted frame-level spectral prediction into the
lowest RTF on both lanes, and RFWave traded iterative solver work for an
intermediate point; the small HiFi-GAN variants minimized footprint below
the leading quality band.  Generation granularity and analytical output
structure tracked the observed orderings better than parameter count
alone, a mechanism-consistent reading, not causal identification, under
native budgets and single trajectories.  The choice of quality proxy is itself a deployment decision: PESQ
and STOI selected RNDVoC while UTMOS selected HiFi-GAN V1 and reshaped
both cost frontiers; proxy agreement still would not constitute listener
preference.  RQ3 inherited this conditionality: the selected artifact
changed with lane and constraint vector.

RQ2 resolved conditionally on operator graph, backend, and
runtime.  FP16 casting halved storage almost uniformly, consistent with its
arithmetic-format origin \citep{micikevicius2018mixedprecision}, yet its
latency effect spanned $0.554\times$ to $1.955\times$; compilation and ONNX
gains depended on graph lowering and kernel choice
\citep{ansel2024pytorch2}; and INT8 coverage was necessary but
insufficient: transformed paths could shrink bytes yet regress speed
\citep{jacob2018quantization}.  Static ONNX INT8 also exposed calibration
dependence: execution-seeded recalibration varied the quantized artifact
by 0.14 PESQ for HiFi-GAN V1; deployed integer pipelines should treat
the calibration set as part of the artifact's identity.  With hosts unpaired, large contrasts are more
credible than marginal crossings, and the taxonomy remains an exploratory
summary of continuous effects.

Dense masks and solver reduction manipulated different objects.  Masks
measured tolerance to zeroed weights, not sparse deployment
\citep{han2015pruning}; recovery showed quality is relearnable,
BigVGAN-base regaining 1.17 PESQ, but the missing dense-continuation
control confounds recovery with added training.  RFWave's step ladder
changed the iteration count and produced the matrix's only smooth,
monotone speed-quality curve, making eight steps a defensible choice.  VocosFormer supports a bounded mechanism hypothesis:
unmasked global attention can spread padded-frame contamination, matching
the padded-to-true-length gap; true-length evaluation removed most of
the quality gap, none of the latency penalty, and attribution requires
masked- and local-attention controls.

In this cohort, post hoc compression split sharply: FP16 held PESQ at
half size in eight of ten groups, while dense masking harmed 18 of 19
groups and static ONNX INT8 three of four; the larger efficiency
separations were architectural, consistent with the architecture-first
emphasis of recent designs
\citep{siuzdak2024vocos,liu2025rfwave,lee2025periodwave,oh2026comvo}.

All intervals are fixed-set descriptive quantities conditional on the
adaptive single-speaker suffix, not an untouched test set; no
multiplicity correction, training replication, listening study, or
external evaluation exists, and warm timing omits streaming, queuing,
cold start, and host identity.  The priority follow-up is an untouched
multi-speaker test set with blinded listening and replicated
trajectories, then randomized host-paired timing with retained raw
repetitions, with sparse kernels, quantization-aware training,
pinned-calibration INT8, and attention masking in those controls.

\section{Conclusion}
\label{sec:conclusion}

Under one corpus, one evaluator, and one measurement protocol, no universal
winner emerged among the twelve finite-budget states (RQ1): RNDVoC
led PESQ and STOI, HiFi-GAN V1 led UTMOS, Vocos led RTF on both lanes,
HiFi-GAN V2 led footprint, and RFWave held the intermediate point.  Deployment transformations were conditional on architecture, operator
graph, and backend (RQ2): FP16 casting was most consistently beneficial,
quantization and compilation were heterogeneous, static ONNX INT8 coupled
large size reductions to material quality loss, dense masks measured
tolerance to zeroed weights rather than sparse execution, and RFWave
solver reduction gave the one monotone speed-quality control.  At the moderate quality, parameter, and storage
vector the artifact-level selector (RQ3) chose FP16-cast Vocos on B200
and dynamic-INT8 Vocos on CPU; the tighter vector selected untransformed
Vocos on CPU, and RNDVoC stayed the highest-PESQ feasible state at both
vectors, never the RTF-minimizing choice.  These conclusions are bounded
by proxy metrics, an adaptive single-speaker evaluation set, single
trajectories, and non-host-paired timing.

\paragraph{Availability.}
The source code of both software packages, the complete experimental
log, and the admission, exclusion, and result registers under
\texttt{artifacts/} are public at
\url{https://github.com/rahul1-bot/VOCODE-Efficient-Neural-Vocoders}; the
graded state is the latest commit on or before 31 July 2026.  LJSpeech is
public under its released license, and resolved run configurations record
the environment.


{\footnotesize
\setlength{\bibsep}{0pt}
\bibliography{assets/research_papers/references}
}

\end{document}
